% language=us runpath=texruns:manuals/workflows

\environment workflows-style

\startcomponent workflows-archives

\startchapter[title=Archives]

\startsection[title=Compression]

The \CONTEXT\ distribution comes with quite some tools and among them are scripts
that can decompress archives. When we started with \LUATEX\ we played with the
idea to run from a zip file and actually that worked so some day we might pick up
on that. Here are a few commands that are available:

\starttyping
mtxrun --script unzip --list       cont-tmf.zip
mtxrun --script unzip --extract    cont-tmf.zip
mtxrun --script unzip --decompress cont-en.fmt
mtxrun --script unzip --decompress foo.gz
mtxrun --script unzip --decompress foo.xz
\stoptyping

The (de)flate and xz decompression are available in the engine and don't depend
on libraries or external programs. We can also work with tarballs:

\starttyping
mtxrun --script untar --list    cont-tmf.tar
mtxrun --script untar --extract cont-tmf.tar.xz
\stoptyping

Both these scripts are reasonable fast and often not much slower than the
independent programs, sometimes even faster.

\stopsection

\startsection[title=\CTAN]

The \CONTEXT\ distributions comes with plenty fonts but you might want to see
what more is available.

\starttyping
mtxrun --script ctan --packages
mtxrun --script ctan --packages --pattern=gyre
\stoptyping

The problem is that there is (end 2026) no way to locate where the files are or
to download a package without maintaining some some registry since we had this
script so we've given up on that and looked at a way to get the files using the
\TEX live repository that has data files at a better defined position. See the
next section for this.

\stopsection

\startsection[title=\TEXLIVE]

The main \TEX\ live installation file has information about packages and when
files come with it. Unfortunately the dependencies or history has resulted in
mixtures of files that are not suitable for \CONTEXT\ so we need to be selective.
Because we're only interested in fonts, we can filter the relevant files (\type
{otf} and \type {ttf}).

\starttyping
mtxrun --script ctan --texlive
\stoptyping

Lists all packages. Based on the name you can decide to focus on one or a few.
You need to be patient because (rather large) compressed files are downloaded
that need to be compressed.

\starttyping
mtxrun --script ctan --packages
mtxrun --script ctan --packages --pattern=pennstander
\stoptyping

You will get a list that indicates what files are available with an indication of
them being already installed or not. With \type {--download} the not yet
installed ones will end up on the current directory, with \type {--force} you get
them all. When there are more packages involved, \type {--all} can be used. When you
install fonts, you best can make a directory, like:

\starttyping
<texroot>/tex/texmf-fonts/fonts/data/somefont
\stoptyping

go there and install in that location. Don't forget to run:

\starttyping
mtxrun --generate
\stoptyping

and if needed:

\starttyping
mtxrun --script font --reload
\stoptyping

Make sure not to install fonts that come with the \CONTEXT\ installer, unless you
know what you're doing.

{\em I might improve this feature over time.}

\stopsection

\stopchapter

\stopcomponent
